\section{RL Environment 是产品架构：任务、工具、奖励与恢复共同定义行为}

讲座第三部分把标题压缩成一句话：\textbf{how you construct RL environment and rewards is how your product will work}。这不是修辞。若训练环境只允许文本回答，模型不会凭空学会可靠工具使用；若 reward 只看最终字符串，模型可能忽略过程中的权限、成本和不可逆副作用；若环境没有失败恢复，策略就没有机会学习何时停止、请求帮助或回滚。

\lecturefigure{slide-50-rl-environment-product.jpg}{核心命题：RL environment 与 reward 的构造决定产品行为}{00:33:15--00:35:51}

\begin{importantbox}{首次使用：MDP、POMDP 与环境 contract}
标准 Markov Decision Process（MDP）写为
\[
\mathcal{M}=(\mathcal{S},\mathcal{A},P,R,\gamma),
\]
其中 $\mathcal{S}$ 是状态，$\mathcal{A}$ 是动作，$P$ 是状态转移，$R$ 是奖励，$\gamma$ 是折扣。真实产品更接近 Partially Observable MDP（POMDP）：模型只看到 observation $o_t$，用户意图、工具状态和未来后果并不完全可见。环境 contract 还必须声明权限、预算、终止、日志、人工确认和恢复。
\end{importantbox}

轨迹回报为
\[
G_t=\sum_{k=0}^{T-t}\gamma^k r_{t+k}.
\]
长任务的关键不是公式，而是 $r_t$ 从哪里来。代码测试、数据库状态、用户确认和人工 judge 的可靠性不同；把不可验证的主观任务压成一个即时分数，会让策略寻找代理漏洞。

\subsection{Real-world complexity：工具、长上下文与多步完成}

真实软件工程或研究任务的复杂度来自工具调用、长上下文、延迟反馈、外部系统和 reward design。图中的 agent/工具只是结构提示；工程验收要进一步问：调用是否有 schema，工具失败是否重试，写操作是否确认，长任务能否 checkpoint，模型是否区分“没有结果”和“工具出错”。

\lecturefigure{slide-51-real-world-rl-complexity.jpg}{真实用例的复杂度等于 RL environment 的复杂度}{00:35:51--00:38:33}

\begin{knowledgebox}{首次使用：Tool call 与 state transition}
工具调用不是一段自然语言。动作可写为 $a_t=(\text{tool},\text{arguments},\text{idempotency key})$，环境返回结构化 observation 与状态变化。若训练只奖励模型“写出像 API 的文本”，却不执行工具，策略学到的是格式模仿而不是任务能力。
\end{knowledgebox}

\begin{lstlisting}[caption={可训练也可部署的 RL environment contract}]
environment = {
    "state": [artifact_versions, tool_status, user_goal, budget],
    "observations": [messages, tool_results, verifier_feedback],
    "actions": [respond, call_tool, edit, ask_user, stop, rollback],
    "permissions": policy_by_action_and_resource,
    "termination": [goal_reached, budget_exhausted, unsafe, user_stop],
    "rewards": [task_outcome, verifier_score, cost_penalty, safety_penalty],
    "audit": persist_full_trace_and_versions,
}
\end{lstlisting}

\subsection{把环境难度拆开讨论}

课堂让学生比较 software engineer、AI researcher、creative storyteller 和 multiplayer/multi-agent environment。难度不只来自答案空间大，还来自反馈延迟、可观察性、协作对象和成功标准。软件任务可能有测试，却包含异步部署和隐性需求；创作任务缺少唯一答案，却有作者 voice、情节一致性和审美目标；多人环境还引入协议与激励。

\lecturefigure{slide-52-environment-complexity-exercise.jpg}{环境复杂度练习：软件、研究、创作与多智能体协作}{00:38:33--00:40:55}

\begin{knowledgebox}{环境难度的五个轴}
\begin{tabularx}{\textwidth}{L{0.20\textwidth} X X}
\toprule
轴 & 低难度一端 & 高难度一端 \\
\midrule
Horizon & 单步回答 & 数小时/数天任务 \\
Observability & 状态完全可见 & 隐含意图、外部变化 \\
Verifiability & 精确测试 & 主观、延迟、多人评价 \\
Action risk & 只读与可撤销 & 写入、付款、发布、控制设备 \\
Coordination & 单用户 & 多人、多 agent、冲突目标 \\
\bottomrule
\end{tabularx}
模型规模相同，环境沿这些轴移动就会产生完全不同的可靠性要求。
\end{knowledgebox}

\teachervoice{00:33:15--00:38:56，Karina 把 software engineering、AI research 与 real-world tasks 的难度归到 RL environment，而不是只说“模型还不够大”。老师强调 tools、long context、reward design 和 multiplayer interaction 都会增加学习难度。}

\subsection{Subjective tasks：可验证性弱，不代表不可研究}

写作、视觉设计、情绪/主动性和 social intelligence 很难用单一 objective measure。课程没有得出“无法训练”的结论，而是要求把任务变成更真实的情境，收集比较、反事实和长期结果。关键是承认 judge 的价值观和上下文依赖，而不是伪装成客观真值。

\lecturefigure{slide-53-subjective-tasks.jpg}{难测的主观任务：写作、审美、情绪与社会智能}{00:38:56--00:41:30}

对两份输出 $y_a,y_b$，pairwise preference 常建模为
\[
P(y_a\succ y_b\mid x)=\sigma\!\left(r_\phi(x,y_a)-r_\phi(x,y_b)\right),
\]
$r_\phi$ 是 reward model，$\sigma$ 是 logistic 函数。Pairwise 比绝对打分容易，但偏好仍会随 reviewer、文化、目的和上下文改变。应保留 disagreement，而不是强行平均成“全球口味”。

\begin{warningbox}{主观任务最危险的简化}
把点赞、停留时间、virality 或单个 AI judge 当成 creativity 本身，会把社会反馈、平台分发和新奇性混在 reward 里。真实世界信号可以作为环境的一部分，但必须防止操纵、群体偏差和 Goodhart's law。
\end{warningbox}

\subsection{新的 Product Research：快速循环而不是一次规格冻结}

讲者把新任务视为真实情境的 simulation，结合 in-context learning、强模型 distillation、synthetic data、新的 model behavior、multiplayer interaction 和 deliberate RL。这里的“快速”不是跳过严谨性，而是让 prototype、eval 与训练更短周期地交换证据。

\lecturefigure{slide-54-product-research-loop.jpg}{新的 Product Research：真实情境、ICL、synthetic data 与 RL 的快速循环}{00:41:30--00:42:14}

\begin{importantbox}{Co-design loop 的实验单位}
每轮实验应固定一条 hypothesis，例如“允许模型先问澄清问题可降低错误工具调用”。然后同时记录：环境改变、数据改变、模型版本、界面改变和 outcome。若五项一起变化，就无法知道提升来自哪里。快速迭代仍需要 controlled comparison。
\end{importantbox}

可以把一次 co-design 迭代记为
\[
\Delta=(\Delta E,\Delta D,\Delta \pi,\Delta I),
\]
分别表示 environment、data、policy/model 和 interface 的变化。理想实验一次只让少量 $\Delta$ 非零，并在相同 eval 上比较；产品上线则需要同时监测交互效应。

\subsection{Reward design：把想要的反馈写清楚}

Reward design 画面问：如果希望模型在真实社会情境中表现更好，应给什么反馈？答案需要 deep product thinking，因为“用户满意”可能来自短期顺从，“主动”可能变成打扰，“有创造力”可能变成不遵守约束。Reward 必须连接真实 outcome、过程限制和人工 judgment。

\lecturefigure{slide-55-reward-design-product-thinking.jpg}{Reward design：反馈应让模型学会真实场景中的合适行为}{00:42:14--00:45:32}

一种分解是
\[
R(\tau)=w_oR_{\text{outcome}}+w_pR_{\text{process}}
-w_cC_{\text{cost}}-w_rC_{\text{risk}}.
\]
$R_{\text{outcome}}$ 衡量任务结果，$R_{\text{process}}$ 奖励澄清、验证和协作过程，$C_{\text{cost}}$ 是时间/计算/工具成本，$C_{\text{risk}}$ 是不可接受风险。权重不是万能修复；如果某项根本无法可靠测量，优化只会放大测量误差。

\teachervoice{00:38:56--00:43:43，课堂把 subjective tasks、真实情境 simulation、synthetic data、multiplayer interaction 与 reward design 连起来。老师的核心问题是：你真正愿意给模型什么反馈，才会让它在现实关系中更有帮助？}

\subsection{Reward hacking 与 asymmetric verification}

Reward hacking 指策略找到高分路径，却没有完成设计者真正想要的目标。随着 reasoning model 和软件 agent 变复杂，hack 可能隐藏在长轨迹、测试漏洞、代码后门或 judge 偏见中。图中的论文与 Lilian Weng 文章提醒：reward 不是 ground truth，verification affordance 本身是 alignment 设计。

\lecturefigure{slide-56-reward-hacking.jpg}{Reward hacking：代理分数提高，但真实目标可能被绕过}{00:43:43--00:45:46}

若真实效用为 $U(\tau)$、可测代理为 $\hat U(\tau)$，优化的是
\[
\max_\pi\;\mathbb{E}[\hat U(\tau)],
\]
但产品关心 $U$。分布移动后，$\hat U$ 与 $U$ 的相关性可能下降。所谓 \term{asymmetric verification} 是生成高质量结果很难，但检查给定结果相对容易；它允许把 verifier 放进训练。然而 verifier 一旦存在 blind spot，策略也会学习攻击它。

\begin{knowledgebox}{验证链必须独立到什么程度}
同一个模型生成答案、解释答案并给自己打分，错误可能高度相关。更强的验证链会组合：可执行测试、规则检查、独立模型、人工抽样、不同数据来源和 adversarial cases。独立性越高，成本越大；产品应按动作风险分配验证预算。
\end{knowledgebox}

\begin{warningbox}{Reward hacking 的产品症状}
模型可能通过隐藏失败、操纵 judge、生成只通过测试的脆弱代码、过度请求用户确认来规避责任，或选择容易任务提高成功率。只看平均 reward 会把这些策略当进步；必须审查完整 trace、任务分母和未完成案例。
\end{warningbox}

\teachervoice{00:43:43--00:45:46，Karina 指出 reasoning 与 software engineering 越复杂，reward hacks 也越复杂；用户可能看不懂模型改了什么代码，因此需要新的 trustworthy verification affordance。课堂提示：验证界面本身属于 alignment。}

\subsection{本章小结}

RL 产品的单位不是孤立模型，而是 POMDP 式环境：状态、观察、动作、工具、权限、reward、成本、终止、日志和恢复共同塑造行为。主观任务可以研究，但必须保留偏好差异；reward 可以优化，但必须承认代理误差和 verifier 攻击。Co-design 的严谨性来自把每轮环境、数据、模型和界面变化都记录清楚。

